\section{Appendix B: Extended Battery Evidence}
\label{app:prof-battery-empirical}

This appendix expands the battery witness with the full publication-facing
empirical package: controller rankings, reliability diagnostics, raw-sequence
track results, operational traces, and the failure-analysis surfaces that
support the compact main-text claim.

\subsection{Locked evaluation scope}
The battery evidence used in the professor-facing manuscript comes from the
locked CPSBench package with 168-hour episodes, seed-controlled runs, and fault
families covering dropout, drift, delay, stale sensing, and combined degraded
telemetry.  The main result is not that every baseline is competitive; it is
that the ORIUS repair-aware variants eliminate true-state violations under the
same locked protocol that exposes failures in the non-repair baselines.

\subsection{Controller ranking on the publication leaderboard}
Table~\ref{tab:prof-battery-leaderboard} records the battery leaderboard as
published in the repo.  The zero-violation tier is occupied by the scenario
methods and the ORIUS repair-aware variants, while the deterministic and static
interval baselines are separated by their non-zero violation rates.

\begin{table}[t]
\centering
\small
\caption{Battery leaderboard from the tracked publication artifact.}
\label{tab:prof-battery-leaderboard}
\begin{tabular}{rlrrr}
\toprule
\textbf{Rank} & \textbf{Controller} & \textbf{Viol.\ rate} & \textbf{IR (\%)} & \textbf{Mean PICP$_{90}$} \\
\midrule
1 & scenario\_mpc & 0.000000 & 0.000000 & 0.934722 \\
2 & scenario\_robust & 0.000000 & 0.000000 & 0.975000 \\
3 & aci\_conformal & 0.000000 & 2.083300 & 0.965278 \\
4 & dc3s\_wrapped & 0.000000 & 0.000000 & 0.880556 \\
5 & dc3s\_ftit & 0.000000 & 8.819400 & 0.930555 \\
6 & deterministic\_lp & 0.129861 & 0.000000 & 0.879861 \\
7 & cvar\_interval & 0.191667 & 0.000000 & 0.879861 \\
8 & robust\_fixed\_interval & 0.191667 & 0.000000 & 0.879861 \\
\bottomrule
\end{tabular}
\end{table}

This ranking is useful because it separates two questions that are often
conflated.  Rank is driven by both safety and cost, while the thesis claim is
only about safety under degraded observation.  On that narrower question, the
important fact is that the ORIUS variants stay in the zero-violation tier and
the deterministic/static baselines do not.

\subsection{Reliability-conditioned baseline audit}
Table~\ref{tab:prof-reliability-baselines} expands the same result family with
coverage and interval-width diagnostics.  The ORIUS FTIT variant is the most
conservative but also the most coverage-stable.  The wrapped variant is less
conservative, but it still preserves zero true-state violations.  That tradeoff
is exactly what the professor-facing manuscript emphasizes: safety is
robust, while efficiency is controller-dependent.

\begin{table}[t]
\centering
\small
\caption{Reliability-conditioned battery baselines from the locked audit.}
\label{tab:prof-reliability-baselines}
\begin{tabular}{lrrrrr}
\toprule
\textbf{Method} & \textbf{TSVR} & \textbf{Obs.\ viol.} & \textbf{Coverage} &
\textbf{Width} & \textbf{Repair rate} \\
\midrule
ORIUS FTIT & 0.000000 & 0.000000 & 0.971627 & 3906.547749 & 0.027976 \\
ORIUS Wrapped & 0.000000 & 0.000000 & 0.884127 & 559.920971 & 0.000000 \\
Adaptive Conformal + Repair & 0.000000 & 0.000000 & 0.962500 & 2824.252435 & 0.013889 \\
Deterministic LP & 0.039087 & 0.039087 & 0.883333 & 503.130805 & 0.000000 \\
Robust Fixed Interval & 0.255754 & 0.255754 & 0.883333 & 503.130805 & 0.000000 \\
\bottomrule
\end{tabular}
\end{table}

The appendix table also makes clear why the paper should not oversell the
deep OQE itself.  The deep reliability model is a useful input to the repair
pipeline, but the published summary shows only modest fault separability
(\texttt{fault\_auc} $=0.534$) and no evidence of a dramatic low-reliability
lift.  In other words, the system is not relying on a miraculous classifier; it
is relying on a conservative safety layer that behaves well even when the
quality score is only moderately informative.

\subsection{Deep OQE diagnostic summary}
Table~\ref{tab:prof-deep-oqe-summary} records the strongest high-level deep OQE
diagnostics from the publication package.

\begin{table}[t]
\centering
\small
\caption{Deep OQE diagnostic summary.}
\label{tab:prof-deep-oqe-summary}
\begin{tabular}{lrr}
\toprule
\textbf{Metric} & \textbf{Heuristic} & \textbf{Deep} \\
\midrule
Eval reliability MAE & 0.117620 & 0.162228 \\
Fault AUC & -- & 0.534482 \\
Low-reliability fault lift & -- & 0.000000 \\
\bottomrule
\end{tabular}
\end{table}

\noindent The value of this table is not that the deep model dominates the
heuristic.  It does not.  The value is that ORIUS remains safe without needing a
perfect fault classifier: the quality signal only has to be informative enough
to steer the repair-and-certify layer away from hidden violations.

\subsection{Raw-sequence track benchmark}
The raw-sequence benchmark is an important negative result because it shows how
hard the unengineered sensing problem remains.  Table~\ref{tab:prof-raw-seq}
compares a reliability-conditioned PatchTST model on raw sequence inputs against
a GBM on engineered tabular features.  Across load, wind, and solar targets, the
raw-sequence model is substantially worse in RMSE and MAE, while calibration is
mixed rather than uniformly superior.

\begin{table}[t]
\centering
\small
\caption{Raw-sequence track benchmark from the publication artifact.}
\label{tab:prof-raw-seq}
\begin{tabular}{llrrrr}
\toprule
\textbf{Target} & \textbf{Track} & \textbf{Model} & \textbf{RMSE} & \textbf{MAE} & \textbf{PICP$_{90}$} \\
\midrule
Load & Raw sequence & PatchTST-reliability & 7414.479 & 5204.866 & 0.610958 \\
Load & Engineered tabular & GBM & 305.121 & 187.869 & 0.882143 \\
Wind & Raw sequence & PatchTST-reliability & 6564.913 & 5091.294 & 0.714088 \\
Wind & Engineered tabular & GBM & 219.108 & 129.362 & 0.936905 \\
Solar & Raw sequence & PatchTST-reliability & 3769.459 & 2533.273 & 0.866943 \\
Solar & Engineered tabular & GBM & 240.353 & 117.319 & 0.787698 \\
\bottomrule
\end{tabular}
\end{table}

\begin{figure}[t]
\centering
\includegraphics[width=\linewidth]{fig_battery_deep_oqe_summary}
\caption{Deep OQE summary plot from the locked publication package.  The main
takeaway is not superiority over the heuristic, but the presence of a usable
quality signal that can be routed into the repair layer.}
\label{fig:prof-deep-oqe-summary}
\end{figure}

\begin{figure}[t]
\centering
\includegraphics[width=\linewidth]{fig_battery_raw_sequence_track_benchmark}
\caption{Raw-sequence track benchmark from the locked publication package.  The
benchmark is intentionally hard: the raw-sequence model is materially worse
than the engineered baseline, which is why ORIUS treats this track as a
stress-test rather than a victory lap.}
\label{fig:prof-raw-seq-benchmark}
\end{figure}

\subsection{Operational trace and failure analysis}
The 48-hour operational trace and the ablation suite give the strongest
mechanistic evidence.  In the trace, interventions cluster around telemetry
degradation rather than random dispatch variation, and the shield is otherwise
transparent.  In the ablation study, disabling OQE, inflation, or projection
reduces the quality of the protection layer; disabling FTIT mostly increases
conservatism rather than breaking safety.

\begin{itemize}
  \item The 48-hour trace records 47 interventions over 1\,680 steps, or 2.8\%
  of steps, with clusters tied to dropout bursts and drift peaks.
  \item The most severe avoided violation in the trace occurs when the LP
  proposes a large discharge against an over-optimistic observed SOC, while the
  shield projects the action back to a safe discharge.
  \item The ablation tables show that OQE and uncertainty inflation are the
  safety-critical elements; FTIT is a performance refinement, not the main
  safety barrier.
\end{itemize}

\noindent This is the right level of detail for the professor-facing manuscript:
the safety story is not a single lucky run, but a reproducible pattern across
controller rankings, reliability audits, raw-sequence stress tests, and trace-level
failure avoidance.
